% ====================================================================
% ch1v22_bib.tex — Chapter 1(흑연+열특성) 장별 참고문헌 (39종) — R1 자동 분할(cite 스캔 기반)
%   원장 = results/V1022_REFERENCE_LEDGER.md (V1 키만). 장 자기완결(리뷰 모음형) — 공통 키의 타 장 중복 수록은 의도(D22).
% ====================================================================
\begin{thebibliography}{99}
\bibitem{dahn1991} J. R. Dahn, ``Phase diagram of Li$_x$C$_6$,'' \emph{Phys. Rev. B} \textbf{44}, 9170 (1991). DOI: 10.1103/PhysRevB.44.9170.
\bibitem{ohzuku1993} T. Ohzuku, Y. Iwakoshi, K. Sawai, ``Formation of Lithium-Graphite Intercalation Compounds in Nonaqueous Electrolytes and Their Application as a Negative Electrode for a Lithium Ion (Shuttlecock) Cell,'' \emph{J. Electrochem. Soc.} \textbf{140}, 2490 (1993). DOI: 10.1149/1.2220849.
\bibitem{bazant2013} M. Z. Bazant, ``Theory of Chemical Kinetics and Charge Transfer based on Nonequilibrium Thermodynamics,'' \emph{Acc. Chem. Res.} \textbf{46}, 1144 (2013). DOI: 10.1021/ar300145c.
\bibitem{eyring1935} H. Eyring, ``The Activated Complex in Chemical Reactions,'' \emph{J. Chem. Phys.} \textbf{3}, 107 (1935). DOI: 10.1063/1.1749604.
\bibitem{glasstone1941} S. Glasstone, K. J. Laidler, H. Eyring, \emph{The Theory of Rate Processes: The Kinetics of Chemical Reactions, Viscosity, Diffusion and Electrochemical Phenomena}, McGraw--Hill (1941).
\bibitem{laidlerking1983} K. J. Laidler, M. C. King, ``The development of transition-state theory,'' \emph{J. Phys. Chem.} \textbf{87}(15), 2657--2664 (1983). DOI: 10.1021/j100238a002.
\bibitem{hill1960} T. L. Hill, \emph{An Introduction to Statistical Thermodynamics}, Addison--Wesley (1960); Dover reprint (1986) --- Ch.~7 (국소화 흡착의 대정준 유도, $q(T)$ 와 Langmuir 등온선).
\bibitem{fowler1939} R. H. Fowler and E. A. Guggenheim, \emph{Statistical Thermodynamics}, Cambridge Univ. Press (1939) --- 국소화 단분자층 흡착의 통계역학 유도.
\bibitem{mcquarrie1976} D. A. McQuarrie, \emph{Statistical Mechanics}, Harper \& Row (1976); Univ. Science Books reprint (2000) --- 독립 부분계 분배함수의 인수분해와 흡착 등온선.
\bibitem{ashcroftmermin1976} N. W. Ashcroft and N. D. Mermin, \emph{Solid State Physics}, Holt, Rinehart and Winston (1976) --- 전자 기체의 Fermi--Dirac 통계와 Sommerfeld 전개(Ch.~2·Appendix C, 장 수준 인용).
\bibitem{mckinnon1983} W. R. McKinnon and R. R. Haering, ``Physical Mechanisms of Intercalation,'' in \emph{Modern Aspects of Electrochemistry}, No.~15 (R. E. White, J. O'M. Bockris, B. E. Conway, eds.), Plenum Press (1983), pp.~235--304 --- 격자기체 삽입 등온선의 원전.
\bibitem{dreyer2010} W. Dreyer \emph{et al.}, ``The thermodynamic origin of hysteresis in insertion batteries,'' \emph{Nature Materials} \textbf{9}, 448 (2010). DOI: 10.1038/nmat2730.
\bibitem{dreyer2011} W. Dreyer, C. Guhlke, and M. Herrmann, ``Hysteresis and phase transition in many-particle storage systems,'' \emph{Contin. Mech. Thermodyn.} \textbf{23}(3), 211--231 (2011). DOI: 10.1007/s00161-010-0178-1. [many-particle 준안정성$\cdot$순차 전환의 이론 상세 --- dreyer2010 의 동반 논문.]
\bibitem{bloom2005} I. Bloom \emph{et al.}, ``Differential voltage analyses of high-power, lithium-ion cells,'' \emph{J. Power Sources} \textbf{139}, 295 (2005). DOI: 10.1016/j.jpowsour.2004.07.021.
\bibitem{weppner_huggins1977} W. Weppner and R. A. Huggins, ``Determination of the Kinetic Parameters of Mixed-Conducting Electrodes and Application to the System Li$_3$Sb,'' \emph{J. Electrochem. Soc.} \textbf{124}(10), 1569--1578 (1977). DOI: 10.1149/1.2133112. [GITT 원전 --- 준평형 OCV 측정법.]
\bibitem{baek_pilon2022} S. W. Baek, M. Saber, A. Van der Ven, and L. Pilon, ``Thermodynamic Analysis and Interpretative Guide to Entropic Potential Measurements of Lithium-Ion Battery Electrodes,'' \emph{J. Phys. Chem. C} \textbf{126}(14), 6096--6110 (2022). DOI: 10.1021/acs.jpcc.1c10414. [전위차 엔트로피 측정$\leftrightarrow$상전이 유형 해석 리뷰, tier B --- load-bearing 주장엔 1차 문헌 병기.]
\bibitem{dubarry2012} M. Dubarry, C. Truchot, B. Y. Liaw, ``Synthesize battery degradation modes via a diagnostic and prognostic model,'' \emph{J. Power Sources} \textbf{219}, 204 (2012). DOI: 10.1016/j.jpowsour.2012.07.016.
\bibitem{swiderska2019} A. {\'S}widerska-Mocek, E. Rudnicka, and A. Lewandowski, ``Temperature coefficients of Li-ion battery single electrode potentials and related entropy changes---revisited,'' \emph{Phys. Chem. Chem. Phys.} \textbf{21}, 2115 (2019). DOI: 10.1039/c8cp06638h. [LiCoO$_2$ 단전극 $\partial\phi/\partial T\approx+0.83$ mV/K, tier B.]
\bibitem{leviaurbach1999} M. D. Levi and D. Aurbach, ``Frumkin intercalation isotherm --- a tool for the description of lithium insertion into host materials: a review,'' \emph{Electrochim. Acta} \textbf{45}, 167--185 (1999). DOI: 10.1016/S0013-4686(99)00202-9. [내재 평형 폭$\cdot$Frumkin 등온선.]
\bibitem{rsc2021} H. Onuma, K. Kubota, S. Muratsubaki, W. Ota, M. Shishkin, H. Sato, K. Yamashita, S. Yasuno, and S. Komaba, ``Phase evolution of electrochemically potassium intercalated graphite,'' \emph{J. Mater. Chem. A} \textbf{9}, 11187--11200 (2021). DOI: 10.1039/D0TA12607A. [저자$\cdot$제목$\cdot$권쪽은 Crossref 로 확정. 흑연 staging 전이별 1차/연속 구분(dilute$\cdot$4L--3L $=$ solid-solution plateau 없음, 나머지 $=$ two-phase plateau)은 이 문헌이 \emph{비교}로 정리한 흑연 삽입 상전이 순서에서 읽음 --- K-graphite 가 주제이나 Li-graphite staging 순서를 대조 제시(tier B, 비교 인용).]
\bibitem{fly2020} A. Fly and R. Chen, ``Rate dependency of incremental capacity analysis (d$Q$/d$V$) as a diagnostic tool for lithium-ion batteries,'' \emph{J. Energy Storage} \textbf{29}, 101329 (2020). DOI: 10.1016/j.est.2020.101329. [C-rate$\uparrow$ 일수록 peak 고전위 이동$\cdot$둔화 --- 동역학 비대칭 꼬리. 저자$\cdot$제목$\cdot$DOI 연도는 Crossref 로 확정(저자 2인 A.~Fly, R.~Chen --- 권 \textbf{29}, 아티클 101329).]
\bibitem{dahn1995} J. R. Dahn, T. Zheng, Y. Liu, and J. S. Xue, ``Mechanisms for Lithium Insertion in Carbonaceous Materials,'' \emph{Science} \textbf{270}(5236), 590--593 (1995). DOI: 10.1126/science.270.5236.590. [탄소 음극의 구조 무질서$\cdot$흑연화도와 리튬 삽입의 \emph{intra-particle} 용량 한계 통계($x_{\max}=1-P$ site-blocking 무질서 모델). \emph{용량-한계 예시}로 인용 --- 비-크기 구조 무질서가 실재함은 보이나, inter-particle 의 apparent-U($\eta$) 분포 형상을 직접 정량하는 1차 근거는 아님(tier C 근거미발견).]
\bibitem{park2021} J. H. Park, H. Yoon, Y. Cho, and C.-Y. Yoo, ``Investigation of Lithium Ion Diffusion of Graphite Anode by the Galvanostatic Intermittent Titration Technique,'' \emph{Materials} \textbf{14}(16), 4683 (2021). DOI: 10.3390/ma14164683. [흑연 하프셀 GITT --- 준평형 OCP 가 입자 크기$\cdot$전극 형상 무관 staging 전위 상수임을 보고; $D_\mathrm{Li}\approx10^{-11}$--$4\times10^{-10}$ cm$^2$/s. 서지는 Crossref 재확인.]
\bibitem{reynier2003} Y. Reynier, R. Yazami, and B. Fultz, ``The entropy and enthalpy of lithium intercalation into graphite,'' \emph{J. Power Sources} \textbf{119--121}, 850--855 (2003). DOI: 10.1016/S0378-7753(03)00285-4. [흑연 삽입 엔트로피$\cdot$엔탈피 실측 --- 부호$\cdot$스케일 anchor(tier B), 전이별 정밀값 아님.]
\bibitem{persson2010} K. Persson, V. A. Sethuraman, L. J. Hardwick, Y. Hinuma, Y. S. Meng, A. van der Ven, V. Srinivasan, R. Kostecki, and G. Ceder, ``Lithium Diffusion in Graphitic Carbon,'' \emph{J. Phys. Chem. Lett.} \textbf{1}, 1176--1180 (2010). DOI: 10.1021/jz100188d. [제일원리$+$실측 Li-흑연 확산 --- 이동 장벽 스케일(tier C).]
\bibitem{persson2010b} K. Persson, Y. Hinuma, Y. S. Meng, A. Van der Ven, and G. Ceder, ``Thermodynamic and kinetic properties of the Li-graphite system from first-principles calculations,'' \emph{Phys. Rev. B} \textbf{82}, 125416 (2010). DOI: 10.1103/PhysRevB.82.125416. [Li-흑연 열역학$\cdot$동역학 제일원리 --- 조성 의존 경향의 근거(tier C).]
\bibitem{cogswell2012} D. A. Cogswell and M. Z. Bazant, ``Coherency Strain and the Kinetics of Phase Separation in LiFePO$_4$ Nanoparticles,'' \emph{ACS Nano} \textbf{6}, 2215--2225 (2012). DOI: 10.1021/nn204177u. [상분리 입자 전극의 계면 에너지 논의 참조 --- 본 문건 $\gamma$ 수치의 출처 아님(스케일 상정은 tier C).]
\bibitem{bernardi1985} D. Bernardi, E. Pawlikowski, J. Newman, ``A General Energy Balance for Battery Systems,'' \emph{J. Electrochem. Soc.} \textbf{132}(1), 5--12 (1985). DOI: 10.1149/1.2113792.
\bibitem{newman} J. Newman, K. E. Thomas-Alyea, \emph{Electrochemical Systems}, 3rd ed., Wiley (2004). (OCV($T$): slope $=\Delta S$, intercept $=\Delta H$; lattice-gas insertion thermodynamics.)
\bibitem{huggins2009} R. A. Huggins, \emph{Advanced Batteries: Materials Science Aspects}, Springer (2009). (격자기체 삽입 전극 열역학·Bragg--Williams 평형 전위.)
\bibitem{allart2018} D. Allart \emph{et al.}, ``Model of Lithium Intercalation into Graphite by Potentiometric Analysis with Equilibrium and Entropy Change Curves of the Graphite Electrode,'' \emph{J. Electrochem. Soc.} \textbf{165}(2), A380 (2018). DOI: 10.1149/2.1251802jes.
\bibitem{occupation2019} M. P. Mercer, M. Otero, M. Ferrer-Huerta, A. Sigal, D. E. Barraco, H. E. Hoster, E. P. M. Leiva, ``Transitions of lithium occupation in graphite: A physically informed model in the dilute lithium occupation limit supported by electrochemical and thermodynamic measurements,'' \emph{Electrochimica Acta} \textbf{324}, 134774 (2019). DOI: 10.1016/j.electacta.2019.134774. [방법 수준 인용 --- 원문 식 미확보.]
\bibitem{chemmater2015} S. Konar, U. Häussermann, G. Svensson, ``Intercalation Compounds from LiH and Graphite: Relative Stability of Metastable Stages\ldots,'' \emph{Chem. Mater.} \textbf{27}(7), 2566--2575 (2015). DOI: 10.1021/acs.chemmater.5b00235. [형성 $\Delta H$ calorimetry --- abstract tier.]
\bibitem{jpcc2021} J. Haruyama, S. Takagi, K. Shimoda, I. Watanabe, K. Sodeyama, T. Ikeshoji, M. Otani, ``Thermodynamic Analysis of Li-Intercalated Graphite by First-Principles Calculations with Vibrational and Configurational Contributions,'' \emph{J. Phys. Chem. C} \textbf{125}(51), 27891--27900 (2021). DOI: 10.1021/acs.jpcc.1c08992. [abstract tier.]
\bibitem{msmr_partI} T. R. Garrick, B. J. Koch, M. Choi, X. Du, A. M. Adeyinka, J. A. Staser, S.-Y. Choe, ``Quantifying the Entropy and Enthalpy of Insertion Materials for Battery Applications Via the Multi-Species, Multi-Reaction Model,'' \emph{J. Electrochem. Soc.} \textbf{171}(2), 023502 (2024). DOI: 10.1149/1945-7111/ad1d27. [★Chapter 1 이 인용하는 ECS Advances 의 ``Part 1''(MCMB 흑연 파라미터화, DOI: 10.1149/2754-2734/ad7d1c)과는 \emph{별개 논문} --- 두 문헌이 공유하는 것은 Part II(DOI: 10.1149/1945-7111/ad70d9) 뿐이다.]
\bibitem{msmr_partII} A. Paul, K. Wolfe, M. W. Verbrugge, B. J. Koch, J. S. Lowe, J. Trembly, J. A. Staser, T. R. Garrick, ``Quantifying the Temperature Dependence of the Multi-Species, Multi-Reaction Model: Part~II. Estimation of Entropy Coefficient for Meso-Carbon Micro-Bead Graphite,'' \emph{J. Electrochem. Soc.} \textbf{171}(10), 103505 (2024). DOI: 10.1149/1945-7111/ad70d9.
\bibitem{standardised2024} A. Hales, J. Bulman, ``A Standardised Potentiometric Method for the Effective Parameterisation of Reversible Heating in a Lithium-Ion Cell,'' \emph{J. Electrochem. Soc.} \textbf{171}(5), 050535 (2024). DOI: 10.1149/1945-7111/ad4918.
\bibitem{hysteresis2018} I. Zilberman, A. Rheinfeld, A. Jossen, ``Uncertainties in entropy due to temperature path dependent voltage hysteresis in Li-ion cells,'' \emph{J. Power Sources} \textbf{395}, 179--184 (2018). DOI: 10.1016/j.jpowsour.2018.05.052.
\bibitem{numverif2026} 본 연구 내부 수치 검증(Derivative A), \code{Anode\_Fit\_v1.0.19} 4-전이 흑연 staging 파라미터, 유한차분 vs.\ 가중식 비교 (2026). [내부 자료 --- 표시 정밀도 일치 확인, 본문 \S\ref{ssec:overlap}$\cdot$\S\ref{ssec:blend} 요지.]
\bibitem{lee2017jcp} K. Lee, S. Lee, C. H. Choi, and S. Lee, ``Effects of external electric field and anisotropic long-range reactivity on charge separation probability,'' \emph{J. Chem. Phys.} \textbf{147}(14), 144111 (2017). DOI: 10.1063/1.5000882. [제2종 Fredholm 방정식의 미지량을 \emph{비}로 치환해 닫는 방법(Eq.~32$\to$34$\to$39)을 적용 --- 부록~\ref{sec:appendix-selfconsistent} 의 ratio 닫힘 원천(적용 논문).]
\bibitem{lee2011jcp} S. Lee, C. Y. Son, J. Sung, and S. Chong, ``Communication: Propagator for diffusive dynamics of an interacting molecular pair,'' \emph{J. Chem. Phys.} \textbf{134}, 121102 (2011). DOI: 10.1063/1.3565476. [제2종 Fredholm 방정식의 효율 해법(ratio 근사) 원 유도이자 \cite{son2013jcp}(Ref.~7)의 선행, \cite{lee2017jcp} Ref.~6 (제목$\cdot$DOI Crossref 확정).]
\bibitem{son2013jcp} C. Y. Son, J. Kim, J.-H. Kim, J. S. Kim, and S. Lee, ``An accurate expression for the rates of diffusion-influenced bimolecular reactions with long-range reactivity,'' \emph{J. Chem. Phys.} \textbf{138}(16), 164123 (2013). DOI: 10.1063/1.4802584. [ratio 닫힘 방법론의 확장 유도 --- \cite{lee2017jcp} Ref.~7 (제목$\cdot$DOI Crossref 확정).]
\bibitem{schmitt2022} J. Schmitt \emph{et al.}, ``Understanding the Influence of Temperature on Phase Evolution during Lithium--Graphite (De-)Intercalation Processes: An Operando X-ray Diffraction Study,'' \emph{ChemElectroChem} \textbf{9}, e202101342 (2022). DOI: 10.1002/celc.202101342. [탈리튬 방향 operando XRD --- stage-2L 형성이 $43\,^\circ$C 뚜렷$\cdot$$0\,^\circ$C 부재(온도 의존 stage-2L 독립 근거). tier B --- 제목$\cdot$저널$\cdot$연도$\cdot$DOI 웹 확인, 저자 전체$\cdot$권쪽은 사용 시 Crossref 최종 대조.]
\end{thebibliography}
% 자산: [B2-154] [B2-155] [B2-156] [B2-157] [B2-158] [B2-159] [B2-160] [B2-161] [B2-162] [B2-163] [B2-164] [B2-165] [B2-166] [B2-167] [B2-168] [B2-169] [B2-170] [B2-171] [B2-172] [B2-173] [B2-174] [B2-175] [B2-176] [B2-177]
% 자산: [C-120] bernardi1985(eq:qrev 원출처) · [C-121] newman(OCV(T) slope=ΔS) · [C-122] huggins2009(격자기체·BW) ·
%   [C-123] bazant2013(regular-solution·Ω=2RT) · [C-124] reynier2003(vib 음baseline) · [C-125] allart2018(정량값·tab:ds) ·
%   [C-126] occupation2019([방법수준]dilute 발산) · [C-127] chemmater2015([형성ΔH abstract]) · [C-128] jpcc2021([abstract]제일원리) ·
%   [C-129] ★msmr_partI(Garrick 2024, Ch1 Part1과 별개·식별함정) · [C-130] msmr_partII(Paul 2024·MCMB) ·
%   [C-131] standardised2024(calorimetry) · [C-132] hysteresis2018(경로의존) · [C-133] ★numverif2026([내부자료]파생A 원천)
